Per a aconseguir soldadures profundes, en la indústria aeroespacial s'utilitza la tècnica de feixos d'electrons d'alta densitat energètica. Aquesta tècnica consisteix a bombardejar amb electrons d'alta energia les peces que s'han de soldar dins d'una cambra de buit. El feix d'electrons es genera escalfant a alta temperatura una banda de wolframi. Posteriorment, el feix s'accelera sota l'acció d'un camp elèctric uniforme que es crea aplicant una diferència de potencial de 15~kV entre l'ànode i el càtode.

\figura[0.5]{fig1.png}

\begin{apartat}{1,25}
Si la separació entre el càtode i l'ànode és d'1,50~cm, determineu el mòdul del camp elèctric que es crea entre l'un i l'altre. Feu un esquema que indiqui la trajectòria dels electrons i la direcció i el sentit del camp elèctric. Quina placa es troba a un potencial més alt, l'ànode o el càtode? Justifiqueu la resposta.
\end{apartat}

\begin{apartat}{1,25}
Considerant que un electró està situat al càtode i parteix del repòs, determineu l'energia i el mòdul de la velocitat de l'electró quan surt de l'ànode. Si la peça que s'ha de soldar es troba al mateix potencial que l'ànode, a quina velocitat impacta l'electró contra aquesta peça?
\end{apartat}

\dades{%
  $m_\mathrm{e} = 9,11\times 10^{-31}\ \mathrm{kg}$.\\
  $|e| = 1,602\times 10^{-19}\ \mathrm{C}$.}
